\scanpage{50}
\mspage{6}
\uline{Remark} In fact, one has a more precise functor, to a category of graphs equipped with a structure in which the $\rho$ part of the weighting is \deleted{\unsure{completed}} \inserted{replaced} by a more precise structure, namely \deleted{a \unsure{map}} \deleted{\inserted{a \unsure{family}}} family
\[
 \deleted{\unsure{\alpha}\,(\widetilde R,\alpha,p:R\longrightarrow\vec A)}
 \qquad (R_a)_{a\in A}
\]
\deleted{\diagram{\begin{tikzcd}[ampersand replacement=\&,row sep=small]\vec A\arrow[d,two heads]\\ \vec A/\sigma\end{tikzcd}}}
of sets, and for $\forall a\in A$, if $\omega(a)$ denotes \unsure{one|the set} of the
\deleted{\unsure{given a permutation of}}
``orientations'' of $a$ ($\operatorname{card}\omega(a)\in\{1,2\}$), a map
\deleted{$R$, and $p$ a map $R\longrightarrow A$ \unsure{conjugate to}}
\[
 \omega\longmapsto u_\omega:\quad \omega(a)\longrightarrow\operatorname{Circ}(R_a)
\]
set of cyclic permutations on $R_a$
\deleted{and \unsure{of another} \unsure{an} \unsure{injective map}}
\[
 \deleted{\unsure{R/\alpha}\lhook\joinrel\longrightarrow\vec A}
\]
with the condition that if $\omega,\omega'\in\omega(a)$, $\omega\ne\omega'$, then
\[
 \deleted{u_\omega u_{\omega'}=\operatorname{id}_R}
 \qquad u_{\omega'}=u_\omega^{-1}.
\]
\margin{$\omega(a)\mathrel{\mathop=^{\rm def}}\{\vec a\in\vec A\mid\text{above }a\}$.}

One can also consider that the datum of the \unsure{supplements} amounts to the datum of
\[
 \begin{array}{rcl}
 r:A&\longrightarrow&\mathbb N\\
 a&\longmapsto&\operatorname{Card}p^{-1}(a)
 \end{array}
\]
(the ``edge-weighting''), and \unsure{the|to the} datum of a family of torsors
\[
 (R_a)_{a\in A},\quad \deleted{\unsure{\omega(a)}}
\]
$R_a$ a torsor under $\mathbb Z/r(a)\mathbb Z\mathbin{\wedge}^{\omega(a)}_{\{\pm1\}}$.
One associates to $(X,\widetilde\Sigma,R)$ such an additional structure:
\[
 R_a=R\cap\widetilde\Sigma_a
\]
($\widetilde\Sigma_a$ the connected component of $\widetilde\Sigma$ \deleted{\unsure{marked}} \inserted{which corresponds to} $a\in A\simeq\pi_0(\widetilde\Sigma)$),
\deleted{$\unsure{R=R(\Sigma C)|R=\pi_0(\widetilde\Sigma)}$, \unsure{$R\longrightarrow S$} \unsure{the datum of} $\widetilde\Sigma$}
\deleted{\diagram{\begin{tikzcd}[ampersand replacement=\&,row sep=small,column sep=large]
R\arrow[r]\arrow[d,equal]\&S\arrow[d,equal]\\
\unsure{\Gamma}\arrow[r]\&\pi_0(\widetilde\Sigma)
\end{tikzcd}}}
$\omega\longmapsto u_\omega$ defined in the obvious way.
\deleted{and $p$ deduced from the cyclic order \unsure{in|on} the $R_a$.}

\uline{Note well} Even if $r(\omega)=1$ or $r(\omega)=2$, the fact of having
\mspage{49}
